\section{Nombres rationnels}
Les paramètres de procédures et les constantes que nous avons utilisés jusqu'à
présent ont des valeurs entières. En \ocaml, on dit qu'ils sont de {\em type}\index{Type}
{\idt int} (de l'anglais \underline{int}eger, entier). Sur les ordinateurs actuels, ces
entiers sont compris entre $-4\,611\,686\,018\,427\,387\,904$ et
$4\,611\,686\,018\,427\,387\,903$, ce qui laisse de la marge pour compter les côtés d'un
polygone ou les carreaux d'une grille. Ils sont codés en base deux dans la mémoire
de l'ordinateur, sur 63 bits.

Nous pouvons également utiliser en \ocaml\ des nombres décimaux avec une virgule comme {\idt 3.14159} appelés flottants (de type \texttt{float}). Ils
s'écrivent selon la convention anglo-saxonne  avec un point à la place de notre
virgule. Les mathématiciens appellent ces nombres des {\em nombres rationnels}\index{Rationnel}. En \ocaml, on
distingue les nombres entiers des nombres rationnels par la présence du point. Par exemple
{\idt 1000} est un entier alors que {\idt 1000.0} est un nombre rationnel, que l'on peut
aussi écrire {\idt 1000.} sans chiffre après le point. Par contre, il faut toujours écrire
au moins un chiffre avant le point~: {\idt .5} est incorrect, il faut écrire {\idt 0.5}.

Les nombres entiers et rationnels sont codés de manières différentes dans la mémoire
de l'ordinateur. Les nombres rationnels sont de type {\idt float}\footnote{\ De
l'anglais {\em floating point number}, nombre à virgule flottante. L'ordinateur
garde en mémoire un nombre fixe de chiffres significatifs et la position de la virgule
(qui ``flotte'' parmi ces chiffres). Ce ne sont pas des nombres
réels puisque les nombres de type {\idt float} ne peuvent pas avoir une infinité de chiffres
différents de zéro après la virgule$\:$!}.

\section{Paramètres de type ``float''}
En \ocaml, il n'est pas nécessaire d'indiquer le type des paramètres d'une
procédure~: le compilateur le devine tout seul (on dit qu'il fait de
l'{\em inférence de types}\index{Inférence de types}) en regardant comment chaque paramètre
est utilisé. Par exemple, dans la procédure~:

\begin{center}
{\mcl let}\ {\idt arc l = lg l; vd ()}
\end{center}

\noindent le paramètre formel {\idt l} (lire la lettre $\ell$ et non le chiffre $1$) est passé à la commande {\idt lg} du \robot\ qui
attend un nombre rationnel. Le compilateur en déduit que {\idt l} est de type {\idt float}.

Le paramètre effectif correspondant doit alors {\em obligatoirement} être un nombre
rationnel, écrit avec un point. Si l'on écrit {\idt arc 100} au lieu de {\idt arc 100.},
le compilateur refuse le programme avec le message
`{\tt The constant 100 has type int but an expression was expected of type float}'
(la constante {\idt 100} est de type {\idt int} alors qu'on attendait une expression de
type {\idt float}) et suggère `{\tt Did you mean 100.?}' (vouliez-vous dire {\idt 100.}$\:$?).
Contrairement à d'autres langages, \ocaml\ ne transforme jamais un entier en rationnel
sans qu'on le lui demande. Pour cela, on utilise la fonction
{\idt float\_of\_int}\index{float\_of\_int}~: si {\idt n} est un entier,
{\idt float\_of\_int n} est le nombre rationnel correspondant (par exemple
{\idt float\_of\_int 3} vaut {\idt 3.}).

Les procédures de commande du \robot\ {\idt lg}, {\idt lgX} et {\idt lgY} (ainsi que
{\idt lgRT}, {\idt rt} et {\idt tl} que nous verrons au chapitre suivant) ont des paramètres
formels de type {\idt float}. C'est pourquoi nous avons toujours écrit {\idt lg 20.} avec un
point. Les commandes {\idt avf}, {\idt cf}, {\idt ec}, {\idt deplacer} ou {\idt tracer}
ont au contraire des paramètres entiers.

\begin{exemple}[Spirale arithmétique]{\rm
La spirale ci-dessous est constituée d'arcs de cercles de rayons successifs
$100$$\:$; $100 / 1.5$ = $66.666\ldots$$\:$; $66.666\ldots / 1.5$ = $44.444\ldots$$\:$;
$44.444\ldots / 1.5$ = $29.629\ldots$$\:$; etc~:

 {\raggedright
 \hbox to 0pt{\hspace*{2cm}\vbox to 0pt{\includevspacePostScript[4cm]{SpiraleArithmetique}{0.8}{0.8}}}
 \input ocaml/SpiraleArithmetique.tex
 }
\noindent $\Box$ }\end{exemple}

\section{Expressions rationnelles}\index{Expression rationnelle (réelle)}
L'addition, la soustraction, la multiplication et la division de nombres rationnels se
notent {\idt +.}, {\idt -.}, {\idt *.} et {\idt /.}, c'est-à-dire comme pour les entiers mais
avec un point à la suite. Par exemple
{\idt (((2.\ *.\ (10.3 +.\ 0.7)) -.\ 1.0) /.\ 3.)} vaut {\idt 7.}

Pourquoi ce point$\:$? En \ocaml, chaque opération a un type bien précis~: {\idt +}
additionne deux entiers et donne un entier, alors que {\idt +.} additionne deux rationnels et
donne un rationnel. Comme les entiers et les rationnels sont codés différemment dans la mémoire
de l'ordinateur, il ne s'agit pas du même calcul. En particulier {\idt 7 / 2} est la division
entière qui vaut {\idt 3} alors que {\idt 7.\ /.\ 2.} vaut {\idt 3.5}. Grâce au point, le
lecteur du programme sait toujours quel calcul est fait. On ne peut donc pas mélanger les
entiers et les rationnels dans une expression~: {\idt 2 *.\ 3.} est refusé par le
compilateur, il faut écrire {\idt 2.\ *.\ 3.} ou, si {\idt n} est un entier,
{\idt float\_of\_int n *.\ 3.}

Si $r$ est un rationnel, {\idt sqrt $r$} est sa racine carrée $\sqrt{r}$ (de l'anglais
\underline{sq}uare \underline{r}oo\underline{t}, racine carrée),
{\idt $r$ ** 2.} est son carré $r^{2}$ (que l'on peut aussi écrire {\idt $r$ *.\ $r$}),
{\idt abs\_float $r$} est sa valeur \underline{abs}olue (pour un entier, on écrit
{\idt abs}), {\idt truncate $r$} est sa partie entière (de l'anglais
\underline{trunc}ate, tronquer) et {\idt round $r$} est l'entier le plus proche (de
l'anglais \underline{round}, arrondir). Les résultats de {\idt truncate} et de {\idt round}
sont des entiers. Par exemple
{\idt truncate 0.3} = {\idt truncate 0.5} = {\idt truncate 0.9} = {\idt 0} tandis que
{\idt round 0.3} = {\idt 0} et {\idt round 0.5} = {\idt round 0.9} = {\idt 1}.
Comme pour les procédures, l'argument d'une fonction qui est une expression doit être
placé entre parenthèses~: {\idt sqrt (x +.\ 1.)}.

Dans une expression rationnelle, on peut utiliser la longueur horizontale {\idt valLgX ()} et la
longueur verticale {\idt valLgY ()} d'un carreau de la grille de déplacement du \robot. Par exemple
{\idt lgX (valLgX ()\ *.\ 2.)} double la taille horizontale d'un carreau de la grille.

\begin{exemple}[Spirale d'or]{\rm
La spirale d'or ci-dessous est constituée d'arcs de cercles successifs dont les rayons
ont pour rapport le {\em nombre d'or\/}
$\Phi = \frac{1}{2}(1 + \sqrt{5}) \simeq 1.618033989$. Pour les tenants d'esthétique
traditionnelle, la spirale d'or constituerait une figure géométrique très harmonieuse
à l'œil.

{\raggedright
 \vbox to 0pt{\includevspacePostScript[4cm]{SpiraleDor}{0.95}{0.95}}
 \input ocaml/SpiraleDor.tex
 }
\noindent $\Box$ }\end{exemple}

\section{Erreurs d'arrondi}
Pour les nombres réels\index{Réel}, qui ont un nombre infini de chiffres après la virgule, comme
$\frac{1}{3} = 0.3333333333333333\ldots$ ou $\pi = 3.1415926535897932\ldots$
l'ordinateur ne peut garder dans sa mémoire que quelques chiffres après la virgule, ce
qui conduit, comme pour une calculette, à des erreurs d'arrondi dans les calculs. Par
exemple en ma\-thé\-ma\-ti\-ques on a~:
\[\left(\left(\left(\left(\left(\sqrt{\sqrt{\sqrt{\sqrt{\sqrt{2}}}}}\right)^{2}\right)^{2}\right)^{2}\right)^{2}\right)^{2} =
\left(\left(\left(\left(\sqrt{\sqrt{\sqrt{\sqrt{2}}}}\right)^{2}\right)^{2}\right)^{2}\right)^{2} =
\] \[\left(\left(\left(\sqrt{\sqrt{\sqrt{2}}}\right)^{2}\right)^{2}\right)^{2} =
\left(\left(\sqrt{\sqrt{2}}\right)^{2}\right)^{2} =
\left(\sqrt{2}\right)^{2} =
2\]
tandis que l'ordinateur trouve  que~:
\begin{center}
{\idt (((((sqrt (sqrt (sqrt (sqrt (sqrt 2.))))) ** 2.)\hspace*{6em}}\\
{\idt \hspace*{6em}** 2.) ** 2.) ** 2.) ** 2.\ = 1.99999999999999645}
\end{center}

\noindent De même, il trouve que {\idt 0.1 +.\ 0.2} vaut
{\idt 0.30000000000000004} (le programme {\idt ErreursDarrondi} du dossier
{\idt livre/\allowbreak 19\_expressions\_rationnelles} fait ces deux calculs).

\noindent  On retrouve des problèmes d'arrondi similaires avec la grille de dé\-pla\-ce\-ment du \robot. De
manière idéale, les coordonnées cartésiennes du \robot\ devraient être réelles.
L'ordinateur en conserve en mémoire une approximation rationnelle. Sur l'écran elles doivent
être arrondies à l'entier le plus proche pour correspondre à un pixel de l'écran.

\noindent \rule{137mm}{0.1mm} %---------------------------------------------
\begin{exercice}{Écrire des programmes \ocaml\  pour réaliser les motifs de la figure
ci-dessous (seul le motif central qui utilise les fonc\-tions tri\-go\-no\-mé\-tri\-ques in\-tro\-dui\-tes
dans l'addenda du chapitre \ref{ch:rec}  est difficile).
\noindent $\Box$}\end{exercice}
\begin{figure}[htbp]
 \includevspacePostScript[4.6cm]{MotifsDeBeard}{0.95}{0.95}
 \caption{Motifs de Horemis et Beard$\:$; spirale de Baravelle}\label{fig:beard}
\end{figure}
\vfill
\begin{corrige}{
Nous donnons un programme, les autres se trouvant dans le dossier
{\idt livre/\allowbreak 19\_expressions\_rationnelles}~:
\label{prog:CarresImbriquesPeints}
\input ocaml/CarresImbriquesPeints.tex
\noindent $\Box$}\end{corrige}

%%%%%%%%%%%%%%%%%%%%%%%%%%%%%%%%%%%%%%%%%%%%%%%%%%%%%%%%%%%%%%%%%%%%%%%%%%%
\noindent \framebox[\textwidth]{
\begin{minipage}{\myparboxwidth}
\setlength{\parindent}{2em}
\small \sf
\vspace{\abovedisplayskip}
\section*{Addenda}
Les nombres rationnels sont souvent qualifiés de réels par abus de langage. Ils
peuvent se noter en utilisant la {\em notation scientifique}\index{Notation scientifique}, qui
consiste à faire
suivre la partie fractionnaire par un facteur d'échelle formé de la lettre {\idt e} ou
{\idt E} (qui se lit ``dix à la puissance'') suivie d'un exposant.

\noindent Par exemple $-3.0${\idt e}$2$ = $-300.0$,
$3.0${\idt e}$0$ = $3.0$ et $3.0${\idt e}$-2$ = $0.03$. Un nombre rationnel ne contient pas de blanc. Par conséquent
$-3.0\hspace{0.3ex}\blanc\hspace{0.3ex}${\idt e}$\hspace{0.3ex}\blanc\hspace{0.3ex}2$ est erroné. En \ocaml, les
nombres de type {\idt float} sont représentés avec 15 à 16 chiffres significatifs en base 10, le plus petit
rationnel strictement positif (sans perte de chiffres significatifs) étant de l'ordre de
$2.2${\idt e}$-308$ et le plus grand de l'ordre de $1.8${\idt e}$308$. Au-delà, le résultat
d'un calcul vaut {\idt infinity} (l'infini).

\vspace{\abovedisplayskip}
\end{minipage}
}
%%%%%%%%%%%%%%%%%%%%%%%%%%%%%%%%%%%%%%%%%%%%%%%%%%%%%%%%%%%%%%%%%%%%%%%%%%%
%%%%%%%%%%%%%%%%%%%%%%%%%%%%%%%%%%%%%%%%%%%%%%%%%%%%%%%%%%%%%%%%%%%%%%%%%%%
\noindent \framebox[\textwidth]{
\begin{minipage}{\myparboxwidth}
\setlength{\parindent}{2em}
\small \sf
\vspace{\abovedisplayskip}
\smallskip

\index{Expression rationnelle (réelle)}\index{Expression entière}La valeur d'une expression arithmétique
peut être de type {\idt int} ou {\idt float}, selon les règles suivantes~:
\begin{list}{$\scriptscriptstyle\bullet$}{
}
\item les constantes entières (sans point) et les opérateurs $+$, $-$, $\ast$, $/$,
{\mcl mod} et {\idt abs}\footnote{\ Nous utilisons des typographies différentes
pour les {\em mots-clés}\index{Identificateur réservé} (comme {\mcl let}
ou {\mcl mod})
et les {\em identificateurs prédéfinis}\index{Identificateur prédéfini}
(comme {\idt sqrt} ou {\idt abs}). Les identificateurs introduits dans un programme \ocaml\
(comme les noms de constantes, de procédures, etc.) doivent être
différents des mots-clés mais peuvent être identiques aux identificateurs prédéfinis. Cette
pratique est déconseillée car les programmes qui changent le sens habituel des identificateurs
prédéfinis sont illisibles.} s'appliquent à des entiers et donnent un résultat entier
de type {\idt int}, de même que {\idt truncate} et {\idt round} qui
s'appliquent à un rationnel$\:$;
\item les constantes rationnelles (avec un point ou un exposant {\idt e} ou
{\idt E}) et les opérateurs $+.$, $-.$, $\ast.$, $/.$, $\ast\ast$, {\idt sqrt} et
{\idt abs\_float} s'appliquent à des rationnels et donnent un résultat rationnel
de type {\idt float}, de même que {\idt float\_of\_int} qui s'applique à un entier.
\end{list}
Une expression qui mélange les deux, comme {\idt 1 +.\ 2.5}, est refusée par le
compilateur.

\smallskip

En \ocaml, les paramètres effectifs doivent avoir exactement le même type que les paramètres
formels. Le compilateur ne transforme jamais de lui-même un entier en rationnel ou un
rationnel en entier~: c'est au programmeur de le demander avec {\idt float\_of\_int},
{\idt truncate} ou {\idt round}. On peut préciser le type d'un paramètre formel lors de la
déclaration de la procédure, comme dans {\mcl let}\ {\idt arc (l :\ float) = lg l; vd ()},
ce qui n'est pas nécessaire mais peut aider le lecteur du programme.

\vspace{\abovedisplayskip}
\end{minipage}
}
%%%%%%%%%%%%%%%%%%%%%%%%%%%%%%%%%%%%%%%%%%%%%%%%%%%%%%%%%%%%%%%%%%%%%%%%%%%
